\input{../tex-inputs/latex-leseansicht-vorspann}

\section[Gustav Schwarzkopf an Arthur Schnitzler, 29. 8. 1913]{L04261 Gustav Schwarzkopf an Arthur Schnitzler, 29. 8. 1913}
\nopagebreak\mylabel{L04261v}
\rehead{ }\normalsize\beginnumbering\briefempfaengerindex{Schnitzler, Arthur@\textsc{Schnitzler, Arthur}!zzzSchwarzkopf, Gustav@\emph{von Gustav Schwarzkopf}!1913-08-291@{29. 8. 1913}|(be}
\toendnotes[C]{\smallbreak\pagebreak[2]}
\correspDesc{Versand  durch Gustav Schwarzkopf am 29. 8. 1913 in Wien
\newline{}Übermittlung  am 29. 8. 1913 in Wien
\newline{}Erhalt  durch Arthur Schnitzler im Zeitraum [30. 8. 1913 – 3. 9. 1913?] in St. Moritz}\toendnotes[C]{\smallbreak}
\Standort{CUL, Schnitzler, B 96a.}
\physDesc{Postkarte, 837 Zeichen
\newline{}Handschrift: schwarze Tinte, deutsche Kurrent
\newline{}Versand: Stempel: »\nobreak{}\oindex{I., Innere Stadt@\textbf{I., Innere Stadt}, \emph{Verwaltungsgebiet}|pwk}1/1 Wien 8, 29. VIII. 13, 8\nobreak{}«.  
\newline{}Schnitzler: mit Bleistift Vermerk: »\textsc{Schwarzk}{[}opf{]}« }\toendnotes[C]{\smallbreak}\pstart{}{\pb}I. Tiefer Graben 17\oindex{Wien@\textbf{Wien}!I., Innere Stadt@\textbf{I., Innere Stadt}!Tiefer Graben 17@\textbf{Tiefer Graben 17}, \emph{Wohngebäude}|pw}.\pend{}\pstart{}II 6.\pend{}{\bigskip}\pstart{}Herrn\pend{}\pstart{}\textsc{D\textsuperscript{r} Arthur Schnitzler}\pend{}\pstart{}\textsc{St. Moriz\oindex{St. Moritz@\textbf{St. Moritz}|pw}}\pend{}\pstart{}(\textsc{Schweiz\oindex{Schweiz@\textbf{Schweiz}|pw}})\pend{}\pstart{}\textsc{Grand Hôtel\oindex{Grand Hotel St. Moritz@\textbf{Grand Hotel St. Moritz}, \emph{Hotel}|pw}}\pend{}{\bigskip}\vspace{1em}
\pstart
           \raggedleft{}{\pb}Wien\oindex{Wien@\textbf{Wien}, \emph{Verwaltungsgebiet}|pw}, 29. VIII 13\pend
           
\pstart{}Lieber Arthur,\pend\vspace{0.5em}
\pstart
           ich war So{\geminationn}tag bei den Kindern\pwindex{Cappellini, Lili 13.\,9.\,1909 Wien – 26.\,7.\,1928 Venedig@\textsc{Cappellini, Lili} (13.\,9.\,1909 Wien – 26.\,7.\,1928 Venedig)|pwv}\pwindex{Schnitzler, Heinrich 9.\,8.\,1902 Hinterbrühl – 12.\,7.\,1982 Wien@\textsc{Schnitzler, Heinrich} (9.\,8.\,1902 Hinterbrühl – 12.\,7.\,1982 Wien), \emph{Regisseur, Schauspieler}|pwv}, deren Ausſehen und
               wolbefinden nichts zu wünſchen übrig läßt. Heini\pwindex{Schnitzler, Heinrich 9.\,8.\,1902 Hinterbrühl – 12.\,7.\,1982 Wien@\textsc{Schnitzler, Heinrich} (9.\,8.\,1902 Hinterbrühl – 12.\,7.\,1982 Wien), \emph{Regisseur, Schauspieler}|pw} hat mich in alle Geheimniſſe \uline{ſeiner}
               Eiſenbahn eingeweiht und Lili\pwindex{Cappellini, Lili 13.\,9.\,1909 Wien – 26.\,7.\,1928 Venedig@\textsc{Cappellini, Lili} (13.\,9.\,1909 Wien – 26.\,7.\,1928 Venedig)|pw} hat mit der
               ganzen Überlegenheit einer erfahrenen Lehrerin mir, dem Ungeſchickten, wol zwanzigmal
               gezeigt, \uline{wie} man vom Schemel ſpringt, ohne zu fallen
               – – ohne aber dabei auch nur einen Augenblick meine ſchützende Hand auszulaſſen. Mir
               geht es vorläufig – wieder ganz gut. Von Ihren überzähligen
               Graden kö{\geminationn}ten wir hier wirklich etwas brauchen, de{\geminationn} ſelbſt an den leidlichen Tagen {\pb}iſt es nicht warm. Ich war auch wieder
               einmal in Baden\oindex{Baden bei Wien@\textbf{Baden bei Wien}, \emph{Hauptstadt}|pw} und habe dort den — »Kuhreigen\pwindex{\textcolor{red}{\textsuperscript{XXXX indx1}}!Kuhreigen@\strich\emph{Der Kuhreigen}|pw}\eventindex{Arena in Baden@\textbf{Arena in Baden}!Aufführung von Der Kuhreigen, 26.8.1913@Aufführung von Der Kuhreigen, 26.8.1913|pwv}«
               gehört ohne ernſtlich Schaden zu nehmen.\pend
           
\pstart
           Ihre Frau\pwindex{Schnitzler, Olga 17.\,1.\,1882 Wien – 13.\,1.\,1970 Lugano@\textsc{Schnitzler, Olga} (17.\,1.\,1882 Wien – 13.\,1.\,1970 Lugano), \emph{Schauspielerin, Sängerin}|pwv} und Sie herzlichſt{\\[\baselineskip]}grüßend{\\[\baselineskip]}Ihr{\\[\baselineskip]}\spacefill\mbox{Gustav.}\pend
           \leftskip=0em{}
\pstart
           \noindent{}Viele Grüße an Frl \textsc{Bachrach\pwindex{Bachrach, Stefanie 22.\,5.\,1887 Wien – 16.\,5.\,1917 ebd.@\textsc{Bachrach, Stefanie} (22.\,5.\,1887 Wien – 16.\,5.\,1917 ebd.), \emph{Krankenpflegerin}|pw}}.\pend
           \selectlanguage{ngerman}\endnumbering\briefempfaengerindex{Schnitzler, Arthur@\textsc{Schnitzler, Arthur}!zzzSchwarzkopf, Gustav@\emph{von Gustav Schwarzkopf}!1913-08-291@{29. 8. 1913}|)be}\mylabel{L04261h}
\begin{anhang}
\end{anhang}\newcommand{\dateiname}{L04261}\newcommand{\titel}{Gustav Schwarzkopf an Arthur Schnitzler, 29. 8. 1913}\newcommand{\editorInnen}{Herausgegeben von Jahnke, SelmaMüller, Martin Anton}\input{../tex-inputs/latex-leseansicht-abspann}
